% SOURCE_AUTHORITY: sga5_fr_workpass.tex, lines 6925--6972 (LF-normalized unit).
% SOURCE_SHA256: 791F4EFFC5E02832D5D77ED03518C8156D6F07E4C8238B03545DB93D883FBB28

Demostremos (c'). Observemos ante todo que, si $Y$ designa un objeto de $P$ y $r$ un entero positivo, existe un objeto $T$ de $P$ tal que $Y=T[r]$ y tal que, para todo objeto $X$ de $P$, el homomorfismo de traslación por $r$
\[
\Hom_P(X,T)\longrightarrow\Hom_P(X[r],Y)
\]
sea una biyección. El sistema proyectivo $(T_n,w_n)$ definido del modo siguiente sirve evidentemente:
Si $Y=(Y_n,u_n)_{n\ge 0}$, $T_n$ es el objeto nulo de $\cC$ para $n\ge r-1$ y $Y_{n-r}$ en caso contrario,
\[
w_n=0\quad\text{para } n\ge r-1,\quad w_n=u_{n-r}.
\]

Dicho esto, supongamos dado un diagrama del tipo que aparece a continuación y demostremos que puede completarse como se indica en el enunciado de (c'):
\[
\begin{tikzcd}[column sep=large,row sep=large]
X[r] \arrow[r,"u'"] \arrow[d,"{v_{rX}}"'] & Y'\\
X .
\end{tikzcd}
\]
Para ello, sea $Y$ un sistema proyectivo tal que $Y'=Y[r]$ y que la aplicación canónica
\[
\Hom_P(X,Y)\longrightarrow\Hom_P(X[r],Y')
\]
sea una biyección. Entonces es claro que el diagrama que aparece a continuación es conmutativo, donde $u$ designa el único morfismo de $X$ en $Y$ tal que $u'=u[r]$:
\[
\begin{tikzcd}[column sep=large,row sep=large]
X[r] \arrow[r,"u'"] \arrow[d,"{v_{rX}}"'] & Y'=Y[r] \arrow[d,"{v_{rY}}"]\\
X \arrow[r,"u"'] & Y .
\end{tikzcd}
\]

Demostremos (d'). Consideremos el diagrama que aparece a continuación, supuesto conmutativo, y demostremos que existe un morfismo $w$ de fuente $Y$ tal que $w\circ f=w\circ g$:
\[
\begin{tikzcd}[column sep=large]
X[r] \arrow[r,"{v_{rX}}"] & X \arrow[r,shift left,"f"] \arrow[r,shift right,"g"'] & Y .
\end{tikzcd}
\]
Ante todo, es claro que el diagrama que aparece a continuación es conmutativo:
\[
\begin{tikzcd}[column sep=large,row sep=large]
X[r] \arrow[r,shift left,"{f[r]}"] \arrow[r,shift right,"{g[r]}"'] \arrow[d,"{v_{rX}}"'] & Y[r] \arrow[d,"{v_{rY}}"]\\
X \arrow[r,shift left,"f"] \arrow[r,shift right,"g"'] & Y .
\end{tikzcd}
\]
En particular, $v_{rY}\circ f[r]=v_{rY}\circ g[r]$. Sea ahora $T$ el sistema asociado a $Y$ y $r$ como al comienzo de la demostración de (c'), de modo que, en particular, $Y=T[r]$. Los trasladados por $r$ de los morfismos $v_{rT}\circ f$ y $v_{rT}\circ g$ son iguales, respectivamente, a $v_{rY}\circ f$ y $v_{rY}\circ g[r]$, y por tanto son iguales. Del hecho de que la traslación por $r$ es una inyección sobre los morfismos de $X$ en $T$ se deduce entonces que se verifica la igualdad que aparece a continuación, de donde se sigue la afirmación:
\[
v_{rT}\circ f=v_{rT}\circ g.
\]
\end{prooftext}
